% ======================================================================
% Ontologie der Kontinua -- Core 1.3.3
% Cross-K-Modul: crossk_global_landscape.tex
% Global cross-level landscape
% ======================================================================

\subsubsection{Globale Cross-K-Landschaft}
\label{sec:crossk-global-landscape}

Die Cross-K-Landschaft beschreibt, wie einzelne Kontinua $K_d$ in eine
einheitliche Multi-Level-Struktur eingebettet sind. Anstatt jede Stufe als
isolierte Theorie zu behandeln, organisiert die Ontologie der Kontinua diese als
eine Kette dimensionssteigernder Übergänge:
\[
R_d : K_d \rightarrow K_{d+1},
\]
wobei jedes $R_d$ ein eingeschränkter Übergang ist, der:
\begin{itemize}
  \item die Kernaxiome (Nicht-Leere von $\Omega$, Schwellen, Flüsse, Zyklen)
        bewahrt,
  \item mindestens eine neue Achse einführt $A_{\text{new}} \in M_d \setminus A(K_d)$,
  \item eine dimensionsbezogene Schwelle überschreitet
        $T_d > \Theta_{d,\text{dim}}$,
  \item einen streng reicheren Raum zulässiger Zustände erzeugt
        $\Omega(K_{d+1}) \supset \Omega(K_d)$.
\end{itemize}

Die globale Landschaft ist daher keine flache Liste von Stufen, sondern eine
strukturierte Folge von Phasenübergängen, in denen jedes $K_d$ aus einem
Vorgänger hervorgeht und seine Nachfolger begrenzt.

% ----------------------------------------------------------------------
\subsubsection{1. Betroffene Stufen und ihre Rollen}

In der groben Auflösung kann die Kette in vier Bänder gegliedert werden:
\begin{enumerate}
  \item \textbf{Physikalischer Bereich} ($K_0$--$K_2$): Grundstruktur, Aktion
        und Perkolation der Konnektivität.
  \item \textbf{Materiell-biologischer Bereich} ($K_3$--$K_5$):
        Felder, Chemie, Kompartimente und proto-neuronale Dynamik.
  \item \textbf{Kognitiver und sozialer Bereich} ($K_6$--$K_8$):
        Bedeutung, soziale Systeme und zivilisatorische Infrastrukturen.
  \item \textbf{Theoretischer Bereich} ($K_9$--$K_{11}$):
        Theorien, Meta-Theorien und trans-meta-theoretische Nebenbedingungen.
\end{enumerate}

Jeder Bereich ist intern kohärent, aber auch durch explizite Cross-Level-Operatoren
$R_d$ und durch gemeinsame Meta-Räume $M_n$ verbunden.

% ----------------------------------------------------------------------
\subsubsection{2. Gemeinsame Achsen und Schwellen über die Stufen hinweg}

Entlang der Kette $K_0$--$K_{11}$ wiederholen sich mehrere strukturelle Muster:
\begin{itemize}
  \item \textbf{Achsen $A(K_d)$}: Jede Stufe führt neue Achsen ein, erbt aber
        transformierte Versionen früherer Achsen (z.\,B. Energie, Gradienten,
        Information, Rollen, Symbole).
  \item \textbf{Schwellen $\Theta(K_d)$}: Existenz-, Stabilitäts-, kritische,
        Todes- und dimensionsbezogene Schwellen treten auf allen Stufen auf,
        jeweils spezialisiert auf den lokalen physikalischen, biologischen,
        sozialen oder theoretischen Kontext.
  \item \textbf{Zustandsdomänen $\Omega(K_d)$}: Jedes Kontinuum besitzt einen
        nicht-leeren Bereich zulässiger Zustände, mit einer Grenze
        $\partial\Omega(K_d)$, die durch verletzte Schwellen definiert wird.
  \item \textbf{Zyklen $C(K_d)$}: Auf jeder Stufe existieren minimale
        stabilisierende Zyklen (von triviale physikalischen Zyklen bis zu
        institutionellen und meta-theoretischen Zyklen).
\end{itemize}

Die globale Landschaft ist als Evolution dieser strukturellen Muster bei steigender
Dimension und Komplexität zu verstehen.

% ----------------------------------------------------------------------
\subsubsection{3. Cross-Level-Flüsse, Zyklen und Spannungen}

Cross-Level-Dynamik lässt sich durch drei Objektfamilien zusammenfassen:
\begin{description}
  \item[Flüsse.] Cross-Level-Flüsse $J_{d,d+1}$ ordnen Zustände und Strukturen
        von $K_d$ auf $K_{d+1}$ zu (z.\,B. chemische Gradienten, die $K_4$
        formen, neuronale Muster, die $K_6$ induzieren, soziale Normen, die
        $K_8$ stützen).
  \item[Zyklen.] Einige Zyklen überbrücken explizit mehrere Stufen, wie etwa
        Zivilisations--Theorie-Rückkopplungen ($K_8$--$K_9$--$K_8$) oder
        Meta-Theorie--Praxis-Zyklen ($K_9$--$K_{10}$--$K_8$).
  \item[Spannung.] Cross-Level-Spannung $T_{d,d+1}$ misst die Inkongruenz
        zwischen den Nebenbedingungen von $K_d$ und $K_{d+1}$; zu große Spannung
        kann die Entstehung einer neuen Stufe verhindern oder einen Zusammenbruch
        nach unten auslösen.
\end{description}

Formal kann man einen generischen Cross-Level-Spannungsfunktional schreiben als
\[
T_{d,d+1} = F\bigl(
  A(K_{d+1}) - A(K_d),\,
  P(K_{d+1}) - P(K_d),\,
  J_{d,d+1},\,
  \Theta(K_{d+1}) - \Theta(K_d)
\bigr),
\]
wobei $T_{d,d+1} > \Theta_{d,\text{dim}}$ einen dimensionsbezogenen
Übergang signalisiert.

% ----------------------------------------------------------------------
\subsubsection{4. Geburts- und Todesbedingungen in der globalen Kette}

Die Geburt einer neuen Stufe $K_{d+1}$ wird durch drei Bedingungen bestimmt:
\begin{enumerate}
  \item \textbf{Verfügbarkeit einer neuen Achse:}
        Es existiert $A_{\text{new}} \in M_d$, so dass
        $A_{\text{new}} \notin A(K_d)$.
  \item \textbf{Dimensionsspannung:}
        Strukturelle Spannung übersteigt eine kritische Schwelle,
        $T_d > \Theta_{d,\text{dim}}$, wodurch die alte Konfiguration ohne
        neue Dimension instabil wird.
  \item \textbf{Erweiterbarkeit des Zustandsraums:}
        Der neue Bereich erfüllt $\Omega(K_{d+1}) \neq \varnothing$ und
        $\Omega(K_{d+1}) \supset \Omega(K_d)$.
\end{enumerate}

Der Tod von $K_d$ kann lokal auftreten (Zusammenbruch eines spezifischen
Kontinuums) oder global (Zerstörung der gesamten Klasse), wenn:
\[
\Omega(K_d) = \varnothing
\quad\text{oder}\quad
k(K_d) \rightarrow 0
\quad\text{oder}\quad
\tau_{\text{cycle}}(K_d) \rightarrow \infty.
\]

Ein globaler Zusammenbruch mehrerer benachbarter Stufen kann als Kaskade von
Schwellenverletzungen verstanden werden, die entlang der Kette propagieren.

% ----------------------------------------------------------------------
\subsubsection{5. Kanonische Cross-Level-Ketten}

Mehrere Cross-Level-Ketten sind von besonderer Bedeutung:
\begin{itemize}
  \item \textbf{Physikalisch-chemisch-biologische Kette:}
        $K_1 \rightarrow K_2 \rightarrow K_3 \rightarrow K_4 \rightarrow K_5$
        (von Aktion und Perkolation bis zu Membranen und proto-neuronaler Dynamik).
  \item \textbf{Biologisch-kognitive-soziale Kette:}
        $K_4 \rightarrow K_5 \rightarrow K_6 \rightarrow K_7$
        (von Gradienten und Impulsen zu Kognition und Institutionen).
  \item \textbf{Sozial-zivilisatorische-theoretische Kette:}
        $K_7 \rightarrow K_8 \rightarrow K_9 \rightarrow K_{10} \rightarrow K_{11}$
        (von Normen und Rollen zu zivilisatorischen Technologien und
        Meta-Rahmenwerken).
\end{itemize}

Die detaillierten Cross-K-Module im technischen Atlas konkretisieren diese
kanonischen Ketten, indem sie für jeden benachbarten Übergang explizite
Darstellungen gemeinsamer Achsen, Schwellen, Flüsse, Zyklen und Todesbedingungen
anbieten.

In diesem Sinn ist die globale Cross-K-Landschaft die \emph{Karte}, wie alle
einzelnen Kontinua $K_d$ in eine einzige kohärente, mehrstufige Struktur passen.
